% BHU PYQ -- Transcribed & Corrected
% Paper: HINMD11: साहित्य एवं सिनेमा
% Exam: Under Graduate Multidisciplinary Course Semester I Examination 2024-25 (NEP)
% Category: Multidisciplinary Course (MDC)
% Department: Hindi / Faculty of Arts

\documentclass[11pt,a4paper]{article}
\usepackage[a4paper,top=2.5cm,bottom=2.5cm,left=2.8cm,right=2.8cm,headheight=14pt]{geometry}
\usepackage{amsmath,amssymb}
\usepackage{fontspec}
\setmainfont{Kohinoor Devanagari}
\usepackage{microtype,fancyhdr,enumitem,hyperref,lastpage,parskip}

\hypersetup{
    colorlinks=true,
    urlcolor=black,
    linkcolor=black,
    pdftitle={HINMD11: साहित्य एवं सिनेमा -- BHU Sem I 2024-25 (NEP MDC)}
}

\pagestyle{fancy}
\fancyhf{}
\renewcommand{\headrulewidth}{0.4pt}
\renewcommand{\footrulewidth}{0.4pt}
\fancyhead[L]{\small \textit{Banaras Hindu University}}
\fancyhead[R]{\small \textit{HINMD11: साहित्य एवं सिनेमा (2024-25)}}
\fancyfoot[C]{\small Page \thepage\ of \pageref{LastPage}}

\newcommand{\pts}[1]{\hfill \textit{[#1]}}

\begin{document}

\begin{center}
  {\large\bfseries BANARAS HINDU UNIVERSITY}\\[2pt]
  {\normalsize Under Graduate Multidisciplinary Course Semester I Examination 2024-25 (NEP)}\\[6pt]
  \rule{0.8\linewidth}{0.5pt}\\[6pt]
  {\large\bfseries HINDI (MDC)}\\[4pt]
  {\large\bfseries Paper Code: HINMD11 : साहित्य एवं सिनेमा}\\[4pt]
  {\normalsize Time: 2:30 Hours\hfill Maximum Marks: 70}
\end{center}

\rule{\linewidth}{0.4pt}

\smallskip

\noindent \textit{सूचना--- सभी प्रश्नों के उत्तर दीजिए।}

\bigskip

\begin{center}
  \textbf{खण्ड --- क}
\end{center}
\noindent \textbf{1. निम्नलिखित प्रश्नों में से किन्हीं तीन प्रश्नों के उत्तर दीजिए:} \pts{3 $\times$ 10 = 30}
\begin{enumerate}[label=(\क)]
  \item भारत में हिन्दी सिनेमा के उदय और विकास की संक्षिप्त रूपरेखा का परिचय दीजिए।
  \item साहित्य और सिनेमा का अंतःसंबंध के स्वरूप पर विचार कीजिए।
  \item सिनेमा की सामाजिक भूमिका पर विचार कीजिए।
  \item साहित्यिक कृतियों पर आधारित सिनेमा के निर्माण में राष्ट्रीय फिल्म विकास निगम की भूमिका स्पष्ट कीजिए।
  \item प्रेमचन्द कृत कहानी 'सद्गति' कहानी में जाति भेद या छुआछूत की समस्या देखने को मिलती है अपने शब्दों में विवेचना कीजिए।
  \item फणीश्वरनाथ रेणु कृत 'तीसरी कसम' कहानी की मूल संवेदना स्पष्ट कीजिए।
\end{enumerate}

\bigskip

\begin{center}
  \textbf{खण्ड --- ख}
\end{center}
\noindent \textbf{2. निम्नलिखित प्रश्नों में से किन्हीं चार लघु उत्तरीय प्रश्नों के उत्तर दीजिए:} \pts{4 $\times$ 6 = 24}
\begin{enumerate}[label=(\क)]
  \item स्वतंत्रता के बाद हिन्दी सिनेमा की प्रमुख प्रवृत्तियों पर प्रकाश डालिए।
  \item बोलती फिल्मों का आरम्भ और विकास पर संक्षिप्त विचार कीजिए।
  \item हिन्दी गीतों की सिनेमा में उपस्थिति स्पष्ट कीजिए।
  \item सिनेमा निर्माण की तकनीकी प्रक्रिया का सामान्य परिचय दीजिए।
  \item 'सद्गति' कहानी का कथ्यबोध स्पष्ट कीजिए।
  \item 'चित्रलेखा' उपन्यास की रचनात्मकता पर संक्षिप्त प्रकाश डालिए।
\end{enumerate}

\bigskip

\begin{center}
  \textbf{खण्ड --- ग}
\end{center}
\noindent \textbf{3. निम्नलिखित प्रश्नों में से किन्हीं चार प्रश्नों के उत्तर दीजिए:} \pts{4 $\times$ 4 = 16}
\begin{enumerate}[label=(\क)]
  \item समानांतर सिनेमा
  \item अमृतलाल नागर
  \item राष्ट्रीय फिल्म विकास निगम (एन.एफ.डी.सी.)
  \item साहित्य और सिनेमा में अन्तर
  \item सिनेमा में लेखक की भूमिका
  \item रचनात्मकता से आप क्या समझते हैं?
\end{enumerate}

\bigskip
\begin{center}
  \textbf{***}
\end{center}

\end{document}
